% Meichen Ding — Project 2 ICS (written by the member; keep the wording, edit only
% if the record changes).
\clearpage
\begin{formbox}
\small
\textbf{Individual Contribution Statement --- Project 2 (PDE BVP)}\\[0.3em]
\textbf{Name:} Meichen Ding \hfill \textbf{Student ID:} 2362344\\
\textbf{Role this project:} Mathematical Theory Lead (model, discretisation
analysis, interface conductance, discrete structure, boundary treatment).

\textbf{1. The specific things I did this project:}
\begin{enumerate}[leftmargin=1.5em,itemsep=1pt,topsep=2pt]
    \item In \S1.2 I fixed the continuous model: the divergence-form problem
    with $0<a_{\min}\le a\le a_{\max}$, the assumptions on $a$ and $g$, the
    transmission condition on the interface $\Sigma$, and the argument for why
    the non-divergence form $-a\Delta u-\nabla a\cdot\nabla u$ loses local
    conservation there. I also wrote the weak formulation and the Lax--Milgram
    well-posedness proof (Theorem~\ref{thm:well_posed}), including the
    Poincar\'e constant used in the coercivity estimate.
    \item In \S2.1 I derived the conservative five-point scheme from the flux
    balance on a control volume, so that the stencil is written in terms of one
    shared conductance per face; this is the structural reason the operator is
    symmetric even when the arithmetic mean is used. I wrote the index map, the
    block-tridiagonal form, the $\operatorname{nnz}=5n^2-4n$ count, and the
    symmetric Dirichlet elimination together with the hand-computable
    $4\times4$ example used as the assembly test.
    \item In \S2.2 I derived the second-order truncation-error theorem with the
    $\theta$-correction of the face coefficient
    ($\theta=0$ midpoint, $\theta=1$ arithmetic, and the perturbation term
    proportional to $a_{xx}u_x$), and the resistor-in-series proposition giving
    the harmonic conductance and the fixed relative error
    $(1-c)^2/(4c)$ of the arithmetic mean. I also wrote the two remarks that
    delimit what these results do and do not test: the smooth benchmark with
    $a=1+x$ is blind to the $\theta$-correction, and the one-dimensional
    resistance chain does not describe the maximum interface-flux error of the
    two-dimensional Dirichlet square.
    \item In \S2.3--\S2.4 I proved that the eliminated matrix is symmetric
    positive definite, an $M$-matrix, and satisfies a discrete maximum
    principle, and I proved the $O(h^2)$ convergence estimate. The maximum-norm
    half of that proof needed a barrier argument rather than an entrywise
    comparison of inverses: I found the counterexample that kills the naive
    argument (Appendix~A, Audit~2) and rewrote the proof with
    $\bm\varphi=\frac12x(1-x)/a_{\min}$ extended by its own boundary values.
    I also derived the spectral bound $\lambda_{\min}\ge 8a_{\min}$ and the
    matched-settings protocol used for the $\star\star$ comparison.
    \item I wrote and ran the independent verification scripts
    \path{verify/check_stability.py}, \path{verify/check_interface_order.py}
    and the one-dimensional reproduction of the layered flux formula, and I
    wrote Appendix~A.
\end{enumerate}

\textbf{2. The hardest conceptual obstacle I encountered:}
The Step ``$A_h\succeq a_{\min}A_h^{(1)}$ and both matrices are $M$-matrices,
hence the inverses compare entrywise'' is the kind of statement that looks like
a one-line consequence and is false. Deciding what was actually true took two
attempts: first a counterexample ($a=1+9\sin^2(\pi x)$ gives an entrywise ratio
of $3.10$), then a different proof route, and finally the observation that the
constant $1/(8a_{\min})$ is attained in one dimension but not in two, where the
true maximum is $0.0736$. Getting the two statements --- the bound is true, the
easy proof is not --- separated was the real work.

\textbf{3. One piece of AI-generated output my team used, and what I changed:}
An AI-assisted draft claimed that arithmetic averaging produces an
$\mathcal{O}(1)$ error in the interface flux while keeping the first-order
nodal error, without distinguishing the face-conductance error from the error
of the solution. I reproduced the one-dimensional chain in closed form
($h(a_L-a_R)^2/(a_L+a_R)^2$ leading term), then re-assembled the full
two-dimensional problem: the nodal and midline flux errors are indeed first
order, but the \emph{maximum} interface-flux error over the square stays at
$0.662$ under refinement, for a different reason (the exact Dirichlet trace on
$y=0,1$). The blanket $\mathcal{O}(1)$ claim was replaced by this two-layer
statement (Appendix~A, Audit~1).

\textbf{4. Anything I should have asked for help with earlier:}
I should have asked the code team for a two-dimensional number before writing an
asymptotic claim that only held in one dimension, and I should have insisted on
per-member commits from Week~5, because a single shared upload account leaves
no traceable history for the contribution statements.

\textbf{Signature:} \underline{\hspace{4.2cm}} \hfill
\textbf{Date:} \underline{\hspace{2.2cm}}
\end{formbox}
